\section{Motivation and Prior Work}\label{sec:motivation_and_prior_work}
Nugget is designed to address the following limitations of prior sampling techniques:

\begin{itemize}
    \item Samples are expensive to find.
    \item Samples are tied to a single executable binary.
    \item Validating the selected samples is infeasible.
\end{itemize}

In this section, we discuss each of these limitations and how prior work in sampling has or has not addressed them.

\subsection{Samples Are Expensive To Find}\label{sec:finding_samples_is_expensive}

Targeted sampling techniques require analyzing the entire workload execution, which is a time-consuming process.

For example, the SPECspeed suite in the \revision{SPEC~CPU2017} benchmark suite~\cite{spec2017} has an average of 21.8 trillion dynamic instructions per benchmark for floating point benchmarks with reference input size, which is the only input size designed to report time metrics~\cite{spec2017-workload-characterization}.
Functional simulators, such as the gem5 atomic model, are often limited to a simulation rate of approximately 1 Million Instructions Per Second (MIPS)~\cite{sustainable_gem5_simulations}.
Therefore, it takes approximately 250 days of \revision{wall-clock} time to analyze an average SPECspeed benchmark with reference input size.
More recent workloads (e.g., the forthcoming SPEC~CPUv8~\cite{spec2017v8}, HPC workloads, and machine learning workloads) are significantly larger, making functional simulation infeasible.

For multi-threaded benchmarks, the analysis overhead is even worse.
We must analyze the program once for every different thread configuration because the number of threads affects the program's execution.
Also, if the workload undergoes any changes (e.g., new compiler, different optimizations, etc.), we must redo the analysis.

Live sampling tools like Pac-Sim~\cite{pacsim} parallelize the phase analysis and detailed simulation to reduce the overall phase analysis penalty.
However, live sampling still relies on the simulator to analyze the program's execution so this technique is constrained by the simulation speed.
If a benchmark takes the simulator 250 days to analyze, as mentioned in the above example, then parallelizing the analysis and detailed simulation will still take at least 250 days.

Tools like DynamoRIO~\cite{dynamorio} and the Intel PinPoints toolkit~\cite{pinpoints} can analyze the program's execution with dynamic binary instrumentation or binary translation.
The instrumentation and profiling routines introduce overhead to the program's execution, but this overhead is significantly smaller than that of simulation (e.g., $4\text{-}40\times$ for QEMU with some single-threaded workloads~\cite{dynamicBinaryInstrumentationForBBVGeneration}).
However, these tools often have restrictions on the architectures they can analyze or the host machines.
Developing architectures that need to be analyzed with a wider range of benchmarks, such as RISC-V, often cannot use these tools because they are not supported.

\begin{tcolorbox}
\textbf{Goal 1:}
Analyze the program's execution quickly and without any restrictions on the architecture or host machines.
\end{tcolorbox}

\subsection{Samples Are Tied To A Single Executable Binary}\label{sec:samples_dont_generalize}

Most prior work relies on binary-specific information, such as the program counter, to define intervals, identify program phases, and select samples.
Any binary change, such as recompiling with different ISA extensions or compilers, alters the static instructions and the dynamic instruction stream, invalidating previously selected samples.
As discussed in Section~\ref{sec:finding_samples_is_expensive}, the high overhead of this analysis makes repeating it for every binary revision time-consuming.

% This is a critical limitation given the growing importance of hardware-software co-design as Moore's law slows~\cite{plenty_of_room_at_the_top}.
% Most targeted sampling techniques produce samples tied to a specific binary, preventing reuse across program binaries.
% For example, to compare a program compiled with and without certain ISA extensions (e.g., RISC-V vector extensions), SimPoint would require rerunning the entire analysis on the new binary.
% Additionally, architectural boundaries are increasingly blurred; ARM-based architectures, for instance, now appear in servers and desktops, emphasizing the need for architecture-independent evaluation~\cite{risc_vs_cisc}.

Cross-binary simulation points~\cite{cross_binary_simpoints} represent a notable attempt to decouple samples from a specific binary.
However, they are limited to single-threaded programs and require significant manual effort, including mapping their control flow graphs and recalculating sample weights.

\begin{tcolorbox}
\textbf{Goal 2:}
Automatically generate samples that are independent of the binary.
\end{tcolorbox}

\subsection{Validating The Selected Samples Is Infeasible}\label{sec:validating_the_selected_samples_is_infeasible}

Sampling methodologies aim to predict overall program performance by measuring only a subset of representative samples and extrapolating the full-program behavior using statistical models. 
Most methodologies focus on predicting metrics that \revision{summarize} the entire hardware performance, such as instructions per cycle (IPC) or total runtime.

To evaluate the accuracy of a set of samples selected by a sampling methodology, the predicted performance metrics are compared against ground truth measurements. 
First, the target performance metric is measured from a full execution of the program on the experimental system, serving as the ground truth. 
Next, the same metric is measured for the selected samples, and statistical processing is applied to estimate the full-program value. 
Finally, the predicted value is compared to the ground truth to quantify the accuracy of the selected samples.

\revision{
The effectiveness of a sample selection method depends on benchmark characteristics. }
\shepherd{Differences in control flow phase behavior, working set size and locality, and input sensitivity can cause a method that succeeds on one workload and input pair to misrepresent another~\cite{characterizing_and_comparing_prevailing_simulation_techniques,simpoint,hot-region-spec2017,xalancbmk,survey_phase_classification,eeckhout-microarchitecture-independent-signatures,MAV}.}
\revision{
For example, empirical studies (e.g., MAV~\cite{MAV}) report outliers where SimPoint fails to capture distinct memory access behavior while SimPoint reports average 3\% IPC error with SPEC~CPU2000~\cite{spec2000} programs~\cite{wavelet-based}. 
Hence, representative intervals should be validated.
However, the only way to validate them is to run the entire workload.
}

\shepherd{Prior sampling techniques are typically evaluated on small benchmarks or truncated runs because their validation depends on simulation~\cite{wavelet-based,taskpoint,simpoint,looppoint,smarts,fsa,cross_binary_simpoints,pacsim,barrier_point,simpoint-plus,parrot}.}
\revision{
A missing piece is an agile path from methodology development to real hardware level validation. 
The intent of the sample set is to approximate the behavior of the real system, so validation by simulation is only meaningful if the simulator faithfully models the target. 
If it does not, the validation does not transfer.
This also leads to difficulty in comparing different methodologies as results are tied to different simulators, forcing redundant engineering effort before we can make apples-to-apples comparisons.
}
For example, LoopPoint is validated with \revision{SPEC~CPU2017} benchmarks using train-size input with Sniper~\cite{sniper}.  
FSA~\cite{fsa} is validated with the first 30~billion instructions of the \revision{SPEC~CPU2006~\cite{spec2006}} benchmarks with gem5~\cite{11_gem5,20_gem5}.
These techniques can only be validated in a limited way because they require a full detailed simulation from beginning to end of the workload to validate.
Detailed simulation can be 2--10$\times$ slower than functional simulation, which implies that validation of full-sized workloads could take nearly a decade! 

\revision{
Lack of an agile workflow to validate quickly has led to sampling techniques being rarely validated with large benchmarks.
}
\shepherd{Evidence suggests that validation results from smaller benchmarks or partial executions do not generalize to larger benchmarks and realistic workloads~\cite{characterizing_and_comparing_prevailing_simulation_techniques,hot-region-spec2017,simpoint,spec2017-workload-characterization,wavelet-based}.}
\revision{
Moreover, the program behavior of a workload with different input can be different~\cite{xalancbmk}.
Thus, per benchmark and input pair validation is needed.
}

\begin{tcolorbox}
\textbf{Goal 3:}
Quickly validating the selected samples for the targeted benchmarks and input size.
\end{tcolorbox}
